\section{Writing logic as constraints}
\label{sec:toolbox}

Many decisions in the case studies are logical: a storage charges or discharges, an appliance
starts or does not, a light is green or red, a task runs on a machine or not. Optimization
solvers do not handle logic; they handle linear inequalities. This chapter collects the
constructions that translate the one into the other. Each uses \emph{binary variables},
$\delta \in \{0,1\}$, and the result is a mixed-integer linear program --- harder than a
continuous one, but far more tractable than a nonlinear one. The constructions are used
repeatedly in \cref{sec:cases}.

\subsection{Indicator variables and the big-M constraint}
\label{sec:bigM}

Let $p \ge 0$ be a continuous variable with a known upper bound $p^{\max}$, and let $\delta$ be
the binary that is meant to indicate whether $p$ is positive. The \emph{indicator constraint}
\begin{equation}
  p \le M\,\delta, \qquad M \ge p^{\max},
  \label{eq:bigM}
\end{equation}
forces $p = 0$ when $\delta = 0$ and leaves $p$ free in $[0, M]$ when $\delta = 1$. It does not
force $\delta = 1$ whenever $p > 0$ by itself; that direction is not needed when the cost makes
an unnecessary $\delta = 1$ never cheaper, which is the usual situation. The constant $M$ should
be as small as the bounds allow, $M = p^{\max}$, so that the relaxation of the problem --- the
one obtained by letting $\delta$ take values in $[0,1]$, which solvers use to guide the search
--- is as tight as possible.

\subsection{Mutually exclusive events}

Two events that cannot happen at the same time --- charging and discharging a battery, importing
and exporting power at a meter --- are written with one indicator each and the constraint
\begin{equation}
  \delta^{a} + \delta^{b} \le 1 .
\end{equation}
Combined with \eqref{eq:bigM} on each power, at most one of the two powers can be positive at a
time. The construction is what makes it possible to write a quantity with two different models
in the two cases, for instance a storage with a charging efficiency and a different discharging
efficiency:
\begin{equation}
\begin{gathered}
  p^{ess} = p^{ch} - p^{dis}, \qquad
  x_{k+1} = x_k + \eta^{ch}\,T\,p^{ch}_k - \tfrac{1}{\eta^{dis}}\,T\,p^{dis}_k, \\
  p^{ch} \le p^{ch,\max}\delta^{ch},\qquad p^{dis} \le p^{dis,\max}\delta^{dis},\qquad
  \delta^{ch} + \delta^{dis} \le 1 ,
\end{gathered}
\end{equation}
which is linear, whereas the same model written with $\max\{0,p\}$ and $\min\{0,p\}$ is not.

\subsection{Semi-continuous variables}

A device that is either off or working between a minimum and a maximum level --- a charger that
cannot deliver an arbitrarily small power --- has a \emph{semi-continuous} variable: $p = 0$ or
$p^{\min} \le p \le p^{\max}$. With one indicator the two inequalities
\begin{equation}
  \delta\,p^{\min} \;\le\; p \;\le\; \delta\,p^{\max}
  \label{eq:semicont}
\end{equation}
force $p = 0$ when $\delta = 0$ and $p \in [p^{\min}, p^{\max}]$ when $\delta = 1$.

\subsection{One start per job: one-hot variables}

A job with a known duration and a known power profile has to be executed once, inside a window
between an earliest start $S$ and a deadline $F$. The decision is \emph{when} to start it, and it
is written with one binary per admissible start time, $s_j \in \{0,1\}$ for $j = S, \dots,
F - N_{job}$, where $N_{job}$ is the duration in steps and the upper limit is the last instant
at which the job still finishes before the deadline. Exactly one of them is one:
\begin{equation}
  \sum_{j=S}^{F-N_{job}} s_j = 1 .
  \label{eq:onehot}
\end{equation}
The power of the job at time $i$ is then the profile placed where the start variable says,
\begin{equation}
  p_i = \sum_{j} s_j\, P(i - j + 1),
  \label{eq:shifted}
\end{equation}
a convolution of the profile $P(\cdot)$ with the one-hot sequence (\cref{fig:onehot}).

\begin{figure}[t]
\centering
\begin{tikzpicture}[x=0.75cm,y=0.9cm,>=Stealth]
  \draw[->] (0,0) -- (14,0) node[right,font=\small] {time};
  \foreach \t in {0,...,13} \draw (\t,0.06) -- (\t,-0.06);
  \node[below,font=\small] at (2,-0.05) {$S$};
  \node[below,font=\small] at (11,-0.05) {$F$};
  \draw[mpc,thick] (2,-0.5) -- (2,2.2);
  \draw[mpc,thick] (11,-0.5) -- (11,2.2);
  \node[mpc,font=\small,above] at (6.5,2.1) {admissible starts $j = S,\dots,F-N_{job}$};
  % one-hot row
  \foreach \j/\v in {2/0,3/0,4/0,5/1,6/0,7/0,8/0} \node[font=\scriptsize] at (\j+0.5,-0.55) {$s_{\j}{=}\v$};
  % profile placed at j=5
  \fill[sun!60] (5,0) rectangle (6,1.4);
  \fill[sun!60] (6,0) rectangle (7,0.9);
  \fill[sun!60] (7,0) rectangle (8,1.1);
  \node[sun!70!black,font=\small] at (6.5,1.75) {$P(1),P(2),P(3)$ placed at $j=5$};
  \draw[decorate,decoration={brace,amplitude=4pt,mirror}] (5,-0.95) -- (8,-0.95) node[midway,below=4pt,font=\small] {$N_{job}=3$ steps};
\end{tikzpicture}
\caption{Load shifting with one-hot start variables. Exactly one $s_j$ equals one inside the
admissible range, and the power profile of the job is placed at that start.}
\label{fig:onehot}
\end{figure}

\subsection{The product of a real variable and a binary}
\label{sec:product}

A term $q\,b$, with $q \in [\underline q, \bar q]$ real and $b \in \{0,1\}$, is not linear. It
is replaced by a new variable $p$ and two pairs of inequalities:
\begin{equation}
  \underline q\, b \le p \le \bar q\, b, \qquad
  q - M(1-b) \le p \le q + M(1-b), \qquad M \ge \max\{\abs{\underline q}, \abs{\bar q}\} .
  \label{eq:product}
\end{equation}
When $b = 0$ the first pair forces $p = 0$ and the second says nothing; when $b = 1$ the second
pair forces $p = q$ and the first is vacuous. The construction recurs in every formulation with
timers, such as the duration of a traffic-light colour, $t_{i+1} = (t_i + T)\,y_i$.

\subsection{Maxima and absolute values}

A goal of the form ``minimize the largest of several quantities'' is not a sum. It is handled
with an \emph{epigraph} variable: to minimize $\max_v w_v$ over a set of quantities $w_v$,
introduce $\bar w$ with the constraints $w_v \le \bar w$ for every $v$, and minimize $\bar w$.
The bound is pushed down until it rests on the largest $w_v$. The same device writes an
absolute value in a cost: $\abs{x}$ is replaced by $x^{+} + x^{-}$ with $x = x^{+} - x^{-}$,
$x^{+}, x^{-} \ge 0$, exact as long as the cost is increasing in both; this is how a tracking
error $\abs{P - P^{ref}}$ enters a linear program. A convex piecewise-linear cost, such as an
approximation of a square, is written in the same way with several pieces of increasing price.

\subsection{Soft constraints and slack variables}
\label{sec:soft}

A constraint that may be impossible to satisfy --- a deadline for a charge when the vehicle
arrives too late, a stop requested in the only hours in which a load can run --- must not be
written as a hard constraint, or the problem becomes infeasible and the controller has nothing
to apply. It is \emph{relaxed} with a non-negative slack variable that measures the violation,
and the violation is paid for in the cost:
\begin{equation}
  g(x) \le 0 \quad\longrightarrow\quad g(x) \le \varepsilon,\ \ \varepsilon \ge 0,
  \qquad J \leftarrow J + d\,\varepsilon \ \text{ or } \ J + d\,\varepsilon^2 .
  \label{eq:slack}
\end{equation}
The weight $d$ sets the behaviour: small, and the controller violates the constraint when that
improves the other goals; large, and it violates it only when the constraint cannot be met at
all. With several relaxed constraints the quadratic form $\sum_p \varepsilon_p^2$ is preferable
to the linear one, because it prefers many small violations to one large one --- it is fairer.

\begin{reading}
\refkey{Hax\_1977}, chapter 9, for integer programming and the modelling of logical conditions:
indicator variables, big-M, mutually exclusive alternatives, fixed charges.
\refkey{Borrelli\_2017} for the same constructions inside an MPC problem, in the framework of
mixed logical dynamical systems. \refkey{Boyd\_2004}, section 4.1.3, for the epigraph form of an
optimization problem.
\end{reading}
